% Generated by tools/edn_latex.py from reports/edn.md.
% Do not edit: change reports/edn.md and run `make edn`.
\ednentry{
  id = {EDN-107},
  title = {The memory alert watches the processes' RSS, at 85\,\% of the limit},
  date = {2026-11-17},
  milestone = {M6: Monitoring},
  topic = {Monitoring, Alerting},
  participants = {Kevin Adameit (decided), AI agents (diagnosis on the VM, options, implementation, review)},
  decision = {ContainerMemoryNearLimit and the panel of each container's share of its limit divide \texttt{container\_\allowbreak{}memory\_\allowbreak{}rss}, the memory the processes allocated themselves, by the container's limit, instead of the working set, and fire at 85\,\% instead of 90\,\%.
The API memory panel draws both, so the kernel's caches stay visible.},
  rationale = {\emph{Alternatives considered:}
\begin{itemize}[leftmargin=*, topsep=2pt]
\item A lighter healthcheck, a raw socket instead of \texttt{urllib.\allowbreak{}request}: it only slows the growth, from about 640 to 240 MiB a day, and changes the code that can kill the server (EDN-93).
\item The working set minus \texttt{container\_\allowbreak{}memory\_\allowbreak{}kernel\_\allowbreak{}usage}: it still counts the log's file cache, and cAdvisor reports that kernel memory as 0 on cgroup v2, so CI would test another figure.
\end{itemize}
\par
\emph{Why:} The alert fired for about 4 of 7 days while the API never grew.
The working set counts kernel caches, and every healthcheck probe, a new Python process every 5 s, leaves about 38 KiB of them: on 2026-11-16 the API held 246 MiB of RSS and 480 MiB of kernel memory.
RSS leaves out the tmpfs, so with its 80 MiB full the API dies from about 88\,\%, hence 85\,\%.
Over those 7 days the API's RSS peaked at 34\,\%.},
  ai = {Information seeking, Alternative generation, Alternative assessment, Recommendation, Solution generation},
  response = {Accepted},
  assessment = {AI first blamed the prediction log's page cache, which its measurement on the VM refuted, then traced the memory to kernel caches and measured each probe's cost.
An independent AI review found the tmpfs gap and a wrong claim about cAdvisor, which led to 85\,\% and a corrected alternative.
Kevin asked for the trade-offs and a test under real conditions, then took both recommendations unchanged.},
  aievidence = {Claude Code session of 2026-11-16 and 2026-11-17.
Question (paraphrased): ``The API's working set fills its limit with kernel cache from the healthcheck's Python processes, not a leak. Which option: alert on RSS instead of the working set, make the healthcheck lighter, or both?''
Kevin answered: ``Base the alert and the panel on the processes' own memory (\texttt{container\_\allowbreak{}memory\_\allowbreak{}rss}) instead of the working set.''},
  evidence = {\href{https://github.com/mlops-2627q1-mds-upc/MLOPS_Recommenditos/issues/368}{Issue \#368}, \href{https://github.com/mlops-2627q1-mds-upc/MLOPS_Recommenditos/pull/369}{PR \#369}, \href{https://github.com/mlops-2627q1-mds-upc/MLOPS_Recommenditos/blob/main/docs/docs/monitoring.md\#containermemorynearlimit}{the alert's runbook}, EDN-93, EDN-99.},
}
